\textbf{Answer to the reproduction question: largely yes.} Secs.~4.4--4.6 of \citet{Kiss2024} are reproduced by an independent Python
implementation: coverage plots, initial costs, Schroeder crest factor and Monte-Carlo trend agree
(T1, T3--T5 reproduced; T6 Schroeder crest factor reproduced, optimised crest factor partially). The converged optimised cost is lower than published (\mCfullMed{} vs.\ 3.88);
budget-limited optimisers stop in the paper's cost range, so T2 and T7 are partially reproduced; early
stopping by the authors (H1) is consistent with this but cannot be proven without their settings. The stated 69$\times$ should read 74$\times$.

\textbf{Extension.} Better coverage does give better models: with the same input spectrum, the model
identified from the space-filling input is \mFOfullforty$\times$ more accurate over the region of
interest (40\,dB, \mFOfullinf$\times$ noise-free) and much more robust to larger inputs, at a moderate
cost (about one third higher free-run error) on test signals of the random-phase kind. The benefit survives when the true system is unknown and the design is computed on an
identified model (\mFOsurforty$\times$, \mFOiterforty$\times$ with one re-identification round), and the
space-filling cost ranks design families by model error ($\rho=\mRhoforty$) but does not resolve
differences among already well-covering designs. The gain survives an actuator peak limit at the Schroeder
crest factor (\mPkOneNineF$\times$), and the remaining noise error is variance that periodic averaging removes,
not a bias that needs instrumental variables.
